% B68 accepted after root independent source, proof, and certificate review.
\subsection*{The surviving full lift is a filling point push}

\statement{B.68}{Proved}{The intermediate relation lifts to a pseudo-Anosov element of the full rigid target}
Let $w=[(ae)^{-1}f(ae),c]$ be the actual cyclic-source word
of B.67, and let $b=\theta_5(w)$.
After moving label $5$ to the last strand, $b$ belongs to
the Birman kernel of $PB_6\to PB_5$.
The corresponding loop in the six-punctured sphere is
not conjugate into a proper free factor of its free
fundamental group $F_5$.  It fills that sphere.
Consequently the image of $b$ in the pure mapping class
group of the seven-punctured sphere is pseudo-Anosov.
This describes the full lift of the additional intermediate
relation; it does not establish injectivity of $\theta_5$
or any assertion about a profinite centralizer.

\begin{proof}
We use the actual common-base words of B.61 and the
strand convention of B.65.  The finite labels in the
six-strand target are $(4,11,5,1,10,6)$.
Conjugating the braid by $\tau=\sigma_3\sigma_4\sigma_5$
in the convention $b\mapsto\tau^{-1}b\tau$ gives the
order $(4,11,1,10,6,5)$.
B.67 proves that deleting its last strand gives the
identity in $PB_5$.  Thus this is a genuine last-strand
point push, before passage to a central quotient.

Here is an explicit convention check for reading its
free-fiber word.  Let $x_1,\ldots,x_6$ be the free
generators used by the faithful Artin action.  Write
the last Artin image as $q x_6 q^{-1}$ and erase
$x_6^{\pm1}$ from $q$; after free reduction call the
result $p\in F(x_1,\ldots,x_5)$.
For elements of the Birman kernel this extraction is
an antihomomorphism: the Artin action induces the
identity after killing $x_6$, so the conjugator formula
for a product reverses the two extracted factors.
For the standard free generators
\[
 A_{i,6}=(\sigma_5\cdots\sigma_{i+1})\sigma_i^2
                  (\sigma_5\cdots\sigma_{i+1})^{-1},
 \qquad 1\leq i\leq5,
\]
it gives
\[
 p(A_{i,6})=(x_{i+1}\cdots x_5)^{-1}x_i
                         (x_{i+1}\cdots x_5).                 \tag{B.68.1}
\]
These five words are a free basis, by a triangular
change of basis, starting with $i=5$ and descending.
Inversion therefore gives an isomorphism from the
Birman kernel to $F_5$.  The actual based loop is
represented by $p^{-1}$ in this convention.
Conjugation and inversion do not affect membership
in a conjugate proper free factor.

The complete finite certificate is
\nolinkurl{next_joint_rigid_symplectic/verify_full_Q_word_filling.py},
its actual output, and
\nolinkurl{full_Q_word_filling_verification.json}.
The verifier recomputes the six-strand braid directly
from the original common-base eleven-strand words;
it does not take the earlier Whitehead search output
as the source of that braid.
It verifies the deletion, the free-fiber word, all five
formulas (B.68.1), and the antihomomorphism convention.
The extracted word has length $144$, with cyclically
reduced length $142$.

For reproducibility, a signed integer $j$ below means
$x_j$ for $j>0$ and $x_{-j}^{-1}$ for $j<0$.
For a signed letter $s$ and a set $A$ with
$s\in A$, $-s\notin A$, use the Whitehead automorphism
that fixes $x_{|s|}$ and sends each other positive
generator $x_j$ to
\[
       s^{-\epsilon_-}\,x_j\,s^{\epsilon_+},\qquad
       \epsilon_-=\mathbf1_{-j\in A},\quad
       \epsilon_+=\mathbf1_{j\in A}.
                                                               \tag{B.68.2}
\]
Here the occurrences of $s$ in (B.68.2) denote the
corresponding signed group letter.  Its inverse has
multiplier $-s$ and set $(A\setminus\{s\})\cup\{-s\}$.
Both inverse compositions are checked on all five
generators at every step.  Applying the following
six automorphisms and cyclically reducing gives:
\[
\begin{array}{c|c|c}
s&A&\text{cyclic length}\\ \hline
3&\{-5,-4,3,4,5\}&98\\
2&\{-5,-4,2,4,5\}&62\\
1&\{-4,1\}&48\\
5&\{-3,-2,2,3,5\}&34\\
4&\{-3,-2,2,4\}&26\\
5&\{-2,5\}&24
\end{array}                                                     \tag{B.68.3}
\]
The final cyclic word is
\[
\begin{split}
 &(4,-1,-4,-3,5,-2,4,2,5,-2,-4,2,\\
 &\hspace{18mm}-5,3,1,-3,-2,5,4,-5,-4,-5,2,3).
\end{split}                                                     \tag{B.68.4}
\]
No minimality assertion about this word is needed.

Form its undirected Whitehead graph on the ten vertices
$\{\pm1,\ldots,\pm5\}$, with an edge from $-v_i$ to
$v_{i+1}$ for each cyclic adjacent pair in (B.68.4).
It has the following fifteen distinct edges:
\[
\begin{array}{rrrrr}
(-5,-4)&(-5,-2)&(-5,4)&(-4,-1)&(-4,1)\\
(-4,2)&(-4,5)&(-3,-1)&(-3,1)&(-3,4)\\
(-2,3)&(-2,5)&(2,4)&(2,5)&(3,5).
\end{array}                                                     \tag{B.68.5}
\]
The graph is connected and remains connected after
deletion of any one of its ten vertices.  The verifier
checks these eleven connectivity statements and saves
a spanning tree for each of them.
Our edge convention is the global sign reversal of
the convention in Heusener and Weidmann, and hence
has the same connectivity properties.

The separable-set form of the cut-vertex lemma states
that a set whose conjugacy classes lie in factors of
a nontrivial free-product decomposition has disconnected
Whitehead graph or a cut vertex
\cite[Corollary~3.5]{DE_HeusenerWeidmann2019}.
Applied to the singleton (B.68.4), its contrapositive
shows that this word lies in no conjugate proper free
factor.  This is the separable-set statement, stronger
than the criterion for a primitive word.
The property is invariant under automorphisms,
conjugation, and inversion, so it also holds for the
actual fiber loop represented by $p^{-1}$.

It remains to connect this algebraic property to
filling on the relevant surface.  The base of this
point push is the sphere with the five finite punctures
$(4,11,1,10,6)$ and the puncture at infinity, label $2$.
Suppose that the loop were disjoint up to homotopy
from an essential simple closed curve $\gamma$.
Both sides of $\gamma$ contain at least two punctures.
After a homeomorphism, the side not containing infinity
encloses the first $k$ finite punctures, with
$2\leq k\leq4$.
Its fundamental group is the proper free factor
$\langle x_1,\ldots,x_k\rangle$ of $F_5$.
The group of the other side is conjugate to
\[
 \langle x_1\cdots x_k,x_{k+1},\ldots,x_5\rangle.
                                                               \tag{B.68.6}
\]
This too is a proper free factor: $x_1\cdots x_k$
is primitive in $\langle x_1,\ldots,x_k\rangle$,
and its rank is $6-k<5$.
A connected representative disjoint from $\gamma$
lies in one of these two sides.  Its conjugacy class
would therefore lie in a proper free factor, contrary
to the certified conclusion above.
The loop consequently intersects every essential
simple closed curve, which is the filling condition.

Kra's point-pushing theorem identifies the point push
of a filling loop on this nonexceptional surface with
a pseudo-Anosov mapping class
\cite[Theorem~2$'$ and Theorem~2(b)]{DE_Kra1981}.
It applies to the six-punctured sphere here and gives
the stated conclusion in the seven-punctured target.
All these statements concern the actual discrete
point push.  They do not identify the kernel of
$\theta_5$, the kernel of its completion, or a
centralizer in a completed mapping class group.
\end{proof}
